\documentclass[11pt]{article}
\usepackage[a4paper,margin=27mm]{geometry}
\usepackage{amsmath,amssymb,booktabs,hyperref}
\title{A braided obstruction and conditional relative-center arithmetic
for $R6_{9,3748}$}
\author{}\date{1 October 2026}
\begin{document}\maketitle

\noindent\textbf{Scope.} The first structural lemma below excludes
\emph{braided} categorifications of $R6_{9,3748}$. The relative-center
calculation that follows is conditional: it does not prove the proposed
categorical reduction and does not give a general exclusion or
categorification of $R6_{9,3748}$.
Of fifteen conditional rank-eighteen based rings, thirteen fail the
formal-codegree $d$-number criterion and two survive. The two survivors
are exchanged by the two tensor factors.

\paragraph{Exact source tensor.}
On labels $[0,1,5]$, the input has the rank-three subring $A$ with basis
$1,a,b$ and
\[
 a^2=1+b,\qquad ab=a+b,\qquad b^2=1+a+b.
\]
The remaining labels $[2,3,4]$ are $X,Xa,Xb$ and form the regular
$A$-module. The full ambient ring, rather than only this local profile,
is checked to be
\[
 R=A[X]/\bigl(X^2-1-(2+2a+b)X\bigr).
\]
All six basis elements are self-dual. The script reconstructs the
complete tensor and compares it under the displayed permutation with
the census tensor. The original tensor is retained in
\texttt{input\_tensor.json}.

\paragraph{Independent obstruction to a braided categorification.}
\textbf{Proposition.} The input based ring has no braided fusion
categorification over an algebraically closed field of characteristic
zero. This conclusion requires none of the relative-center hypotheses
below.

\emph{Proof.} Suppose $\mathcal C$ were such a braided categorification,
and let $\mathcal D$ be its fusion subcategory on labels $[0,1,5]$.
The fusion rules of $A$ show that either nonunit simple generates all
of $\mathcal D$: $a^2$ contains $b$, and $b^2$ contains $a$.
Thus $\mathcal D$ has no proper nontrivial fusion subcategory.
Its M\"uger center is a fusion subcategory, hence is either
$\mathrm{Vec}$ or $\mathcal D$.

The second case would make $\mathcal D$ symmetric. Deligne's theorem
identifies a symmetric fusion category in characteristic zero with
$\operatorname{Rep}(G,z)$; in particular it has integral Frobenius--Perron
dimensions \cite[Theorem 2.43 and Corollary 2.46]{DGNO}.
But $\alpha=\operatorname{FPdim}(a)$ satisfies
$\alpha^3-\alpha^2-2\alpha+1=0$. This monic cubic has no rational
root: the only possible rational roots are $1$ and $-1$, and neither
is a root. Thus $\alpha$ is not an integer, a contradiction.
Consequently the M\"uger center of $\mathcal D$ is trivial, so
$\mathcal D$ is nondegenerate \cite[Theorem 3.4]{DGNO}.

Because the \emph{ambient} category is now assumed braided,
\cite[Theorem 3.13]{DGNO} applies and gives
\[
 \mathcal C\simeq\mathcal D\boxtimes\mathcal D',
 \qquad 6=3\operatorname{rank}(\mathcal D').
\]
Therefore $\mathcal C$ has a rank-two fusion subcategory $\mathcal D'$.
This is ruled out directly by its input tensor. If its nonunit simple
had label $j$, closure would require
$\operatorname{supp}(j^2)\subseteq\{0,j\}$. The actual supports are
\[
\begin{array}{c|ccccc}
j&1&2&3&4&5\\\hline
\operatorname{supp}(j^2)&
\{0,5\}&\{0,2,3,4\}&\{0,2,3,4,5\}&
\{0,1,2,3,4,5\}&\{0,1,5\}.
\end{array}
\]
None has the required containment. This proves the proposition.
\hfill$\square$

The proof does not classify the category $\mathcal D$ as a particular
quantum-group category, and does not need the rank-two classification.
Deligne's integrality consequence and the cited DGNO decomposition
theorem do not assume a pivotal, spherical, pseudo-unitary, or unitary
structure. The conclusion concerns braidings only: it does not alter
the existing general, pivotal, spherical, pseudo-unitary, or unitary
classification fields.

\paragraph{Unproved categorical hypotheses.}
For a hypothetical characteristic-zero fusion categorification
$\mathcal C$, write $\mathcal D$ for its subcategory on $1,a,b$ and
$\mathcal M$ for its complementary $\mathcal D$-bimodule.
The calculation assumes the following relative-center description:
\begin{enumerate}
\item $K_0(Z(\mathcal D))$ is identified as a based ring with
$B=A\otimes A$, and the forgetful map is
$F(u\otimes v)=uv$.
\item $\mathcal M$ is the regular $\mathcal D$-bimodule, with $X$
corresponding to the unit. The verified regular left-module fusion
rules alone are not a proof of this bimodule equivalence.
\item Under the induced identifications, the relative center
$Z_{\mathcal D}(\mathcal C)$ has simple basis $B\sqcup Bt$,
$F(t)=X$, and $B$ is central in its Grothendieck ring.
\end{enumerate}
The arithmetic verifier does not construct the required bimodule
equivalence, half-braidings, or categorical equivalences.

Under these hypotheses, the only nonnegative preimage of the unit
under $F:B\to A$ is the basis unit. The regular-module identification
and duality consequently give $t^*=t$, and the $B$-part of $t^2$ is
$1$. Thus the candidate ring must have the form
\[
 E_C=B[t]/(t^2-1-Ct),\qquad C\in B_{\geq0},\quad F(C)=2+2a+b.
\]
Every such algebra is a rank-eighteen based ring. This last assertion
is also verified directly, including all associativity entries, for
every enumerated candidate.

\paragraph{Complete finite domain.}
Write $e=(1,a,b)$ and $C=\sum_{i,j=0}^2c_{ij}e_i\otimes e_j$.
The exact constraints are
\begin{align*}
 c_{00}+c_{11}+c_{22}&=2,\\
 c_{01}+c_{10}+c_{12}+c_{21}+c_{22}&=2,\\
 c_{02}+c_{20}+c_{11}+c_{12}+c_{21}+c_{22}&=1.
\end{align*}
There are fifteen nonnegative integral solutions. Indeed, in the last
equation exactly one of six coefficients is one. Choosing $c_{11}$
gives three solutions, choosing $c_{22}$ gives two, choosing either
$c_{02}$ or $c_{20}$ gives three each, and choosing either $c_{12}$ or
$c_{21}$ gives two each. The script independently enumerates a complete
rectangular domain using coefficient bounds from $F(C)$.
Duality is componentwise and does not exchange the two tensor factors;
no symmetry constraint $c_{ij}=c_{ji}$ is imposed.

\paragraph{Exact formal-codegree calculation.}
For a character of $A$, let $\alpha$ be the value on $a$. Then
\[
 \alpha^3-\alpha^2-2\alpha+1=0,\qquad
 \chi(b)=\alpha^2-1,\qquad f_A=2\alpha^2+\alpha+1.
\]
For each of the nine characters $\chi$ of $B$, put
$c=\chi(C)$ and let $f_B$ be its formal codegree. Both roots of
$\lambda^2-c\lambda-1=0$ extend $\chi$ to $E_C$. The corresponding
formal codegrees are $f_B(1+\lambda^2)$, so their quadratic polynomial is
\[
 z^2-f_B(4+c^2)z+f_B^2(4+c^2).
\]
The forgetful constraint determines $c$ on only the three diagonal
characters of $A\otimes A$. All six off-diagonal characters are
included by the calculation.

No rank-eighteen character field is constructed. In the regular
representation of $B$, let $H$ be multiplication by $C$ and $K$ its
Casimir matrix. The Casimir matrix of $E_C$ is checked directly to be
\[
 \begin{pmatrix}2K&KH\\KH&K(2I+H^2)\end{pmatrix}.
\]
Its characteristic polynomial is factored exactly over $\mathbb Q$.
As an independent formula check, it is compared with
\[
 \det\bigl(z^2 I-zK(4I+H^2)+K^2(4I+H^2)\bigr)
\]
at nineteen distinct integer values of $z$. Both polynomials have
degree eighteen, so these exact evaluations prove equality.
Every irreducible factor and multiplicity is recorded.

For a monic irreducible factor
$p(z)=z^d+a_1z^{d-1}+\cdots+a_d$, the $d$-number test checks
$a_d^j\mid a_j^d$ for every $1\leq j\leq d$.
Every failing divisibility, including its nonzero remainder, is saved.
Each selected exclusion witness is independently accepted by the
package's existing \texttt{rings.check\_dnumber} function. The
formal-codegree theorem and polynomial criterion are
\cite[Theorem 1.2 and Lemma 2.7(v)]{Ostrik}. This is a general fusion-category
obstruction; no pseudo-unitary hypothesis is used for this arithmetic
step. It excludes the specified rank-eighteen candidate only.

\paragraph{Results.}
Cases are in lexicographic order of
$(c_{00},c_{01},c_{02},c_{10},c_{11},c_{12},c_{20},c_{21},c_{22})$.
All fifteen exact tensors pass the based-ring checks. The following is
the complete labelled domain; it is not a claim that distinct rows
give pairwise nonisomorphic rings. Each tuple is one row of the
coefficient matrix $(c_{ij})$.
\begin{center}
\begin{tabular}{rcccl}\toprule
Case & $(c_{00},c_{01},c_{02})$ & $(c_{10},c_{11},c_{12})$
& $(c_{20},c_{21},c_{22})$ & $d$-number result\\\midrule
1  & $(1,0,0)$ & $(1,0,0)$ & $(0,0,1)$ & excluded\\
2  & $(1,0,0)$ & $(2,1,0)$ & $(0,0,0)$ & excluded\\
3  & $(1,1,0)$ & $(0,0,0)$ & $(0,0,1)$ & excluded\\
4  & $(1,1,0)$ & $(1,1,0)$ & $(0,0,0)$ & excluded\\
5  & $(1,2,0)$ & $(0,1,0)$ & $(0,0,0)$ & excluded\\
6  & $(2,0,0)$ & $(1,0,0)$ & $(0,1,0)$ & excluded\\
7  & $(2,0,0)$ & $(1,0,1)$ & $(0,0,0)$ & excluded\\
8  & $(2,0,0)$ & $(2,0,0)$ & $(1,0,0)$ & survives\\
9  & $(2,0,1)$ & $(2,0,0)$ & $(0,0,0)$ & excluded\\
10 & $(2,1,0)$ & $(0,0,0)$ & $(0,1,0)$ & excluded\\
11 & $(2,1,0)$ & $(0,0,1)$ & $(0,0,0)$ & excluded\\
12 & $(2,1,0)$ & $(1,0,0)$ & $(1,0,0)$ & excluded\\
13 & $(2,1,1)$ & $(1,0,0)$ & $(0,0,0)$ & excluded\\
14 & $(2,2,0)$ & $(0,0,0)$ & $(1,0,0)$ & excluded\\
15 & $(2,2,1)$ & $(0,0,0)$ & $(0,0,0)$ & survives\\\bottomrule
\end{tabular}
\end{center}
Thirteen have verified $d$-number obstructions. Each selected witness
fails already at $j=1$ for a degree-six irreducible factor; all factors
and their other failed divisibilities are retained in the JSON report.
The surviving cases are
\[
 \begin{array}{ccl}
 8 &:& C=(2+2a+b)\otimes1,\\
 15&:& C=1\otimes(2+2a+b).
 \end{array}
\]
They have the same three irreducible degree-six formal-codegree
polynomials, each with multiplicity one; all pass the $d$-number test.
At the level of based rings these two candidates are
$R6_{9,3748}\otimes A$, with the factors exchanged. This does not
establish a tensor-product decomposition of a category.

\paragraph{What the one-sided survivors would give categorically.}
Assume the categorical hypotheses above, including their specified
forgetful functor. For $C=(2+2a+b)\otimes1$, the six simple objects
\[
 (u\otimes1),\quad (u\otimes1)t\qquad(u=1,a,b)
\]
form a fusion subcategory $\mathcal S$ of the actual relative center:
the displayed relation proves closure under products, and duals fix
the basis. The restriction $F|_{\mathcal S}:\mathcal S\to\mathcal C$
takes these simples bijectively to all six simples of $\mathcal C$.
It is fully faithful on simple Hom spaces, hence an equivalence.
A monoidal inverse gives a section
$s:\mathcal C\to Z_{\mathcal D}(\mathcal C)$.

This section provides a coherent central lift of $\mathcal D$ into
$Z(\mathcal C)$, rather than merely objectwise commutation. To see
this in strict notation, write its relative half-braiding as
\[
 \gamma^Y_V:VY\longrightarrow YV,\quad
 V\in\mathcal D,\ Y\in\mathcal C.
\]
For fixed $V$ these maps are natural in $Y$, and the monoidality of
$s$ gives the half-braiding identity in the variable $Y$.
The relative-center identity gives tensor compatibility in $V$.
Thus $V\mapsto(V,\{\gamma^Y_V\}_Y)$ is a central tensor lift.
The case with tensor factors exchanged is the same. If the identification
$Z(\mathcal D)\simeq\mathcal D\boxtimes\mathcal D^{\rm rev}$ is
categorically established, its restriction induces the corresponding
one of the two nondegenerate braidings on $\mathcal D$.

This conclusion still does \emph{not} supply a braiding on all of
$\mathcal C$, or the second half-braiding needed to define a M\"uger
centralizer inside $\mathcal C$. The tensor factorization theorem
\cite[Theorem 3.13]{DGNO} assumes a braided ambient category. A general
central embedding of a modular category into the center of an
unbraided fusion category is insufficient. For example, take the
relative center of a Tambara--Yamagami category for the group of order
two over its pointed subcategory. It has four invertible simples and
two simples of dimension $\sqrt2$ by \cite[Section 4.2]{GNN}, hence
rank six and dimension eight. Its four-simple modular neutral category
$Z(\mathrm{Vec}_{\mathbb Z/2})$ has a central lift to the center of the
relative center. A tensor splitting by this four-simple category would
leave a category of dimension two, hence rank two, and would incorrectly
give rank eight. This example rules out the unrestricted implication;
it does not rule out a further theorem specific to $\mathcal D$ here.

For comparison, \cite[Theorem 1.2]{EM} proves that
$\operatorname{BrPic}(\operatorname{Ad}(A_6))$ is trivial, also for
Galois conjugates. Applying that fact here requires a categorical
identification of $\mathcal D$ and invertibility of the relevant
bimodule. Moreover the original decomposition is not a group grading,
since $X^2$ has a nonzero complementary component. The cited
Brauer--Picard computation alone is not a classification of these
nongraded extensions. No rank-two centralizer, scalar-coefficient
necessity, or general exclusion is claimed.

\paragraph{Reproduction and provenance.}
Run \texttt{verify\_relative\_candidates.py} with a SymPy-enabled
Python interpreter to regenerate \texttt{conditional\_certificate.json}.
Run it with \texttt{--verify} to replay the saved exact results without
writing them. The default driver launches each of the fifteen cases
in a separate subprocess with a hard fifty-five-second deadline;
unfinished jobs cannot be counted as exclusions. These cases required
approximately $0.07$--$0.14$ seconds each of computation in the first
run. The certificate records the source tensor and census hashes, the
verifier and arithmetic-library hashes, complete formal polynomials,
all fifteen coefficient vectors, and the exact witnesses. The replay
uses \texttt{python -O}; no correctness check relies on Python
\texttt{assert} statements. The calculation changes no canonical
classification status or existing certificate.
\begin{thebibliography}{9}
\bibitem{Ostrik} V.~Ostrik,
\emph{On formal codegrees of fusion categories},
Math. Res. Lett. 16 (2009), 895--901,
\href{https://arxiv.org/abs/0810.3242}{arXiv:0810.3242}.
\bibitem{GNN} S.~Gelaki, D.~Naidu, D.~Nikshych,
\emph{Centers of graded fusion categories},
Algebra Number Theory 3 (2009), 959--990,
\href{https://arxiv.org/abs/0905.3117}{arXiv:0905.3117}.
\bibitem{DGNO} V.~Drinfeld, S.~Gelaki, D.~Nikshych, V.~Ostrik,
\emph{On braided fusion categories I},
Selecta Math. 16 (2010), 1--119,
\href{https://arxiv.org/abs/0906.0620}{arXiv:0906.0620}.
\bibitem{EM} C.~Edie-Michell,
\emph{The Brauer--Picard groups of fusion categories coming from the
ADE subfactors},
\href{https://arxiv.org/abs/1709.04721}{arXiv:1709.04721}.
\end{thebibliography}
\end{document}
